% Traduction française COMIMa de l'énoncé original (anglais).
On part de l'état $(a, b)$, où $a, b$ sont des entiers strictement positifs, et on lui applique l'algorithme suivant :
\[ \textbf{tant que } a > 0,\ \textbf{si } a < b \textbf{ alors } (a, b) \leftarrow (2a, b - a) \textbf{ sinon } (a, b) \leftarrow (a - b, 2b). \]
Pour quelles positions de départ l'algorithme s'arrête-t-il ? En combien d'étapes, s'il s'arrête ? Que peut-on dire des périodes et des prépériodes ?

Mêmes questions lorsque $a, b$ sont des réels strictement positifs.
